% Source PDF pages 18--22; standard pages 6--10.
% Continuation of Clause 3.1 definitions.

\newpage

\defentry{available dimensions}{}
\begin{hlalist}
    \item Pertaining to an attribute: the dimensions associated with the class attribute in the federation object model (FOM). The available dimensions of an instance attribute are the available dimensions of the corresponding class attribute.
    \item Pertaining to an interaction class: the dimensions associated with that interaction class in the FOM. The available dimensions of a sent interaction are the available dimensions of the interaction class specified in the \hlaService{Send Interaction With Regions} service invocation.
\end{hlalist}
\hlaNote{See 9.1 in IEEE Std 1516.1-2010.}

\defentry{available parameters}{The set of declared parameters of an interaction class in union with the set of inherited parameters of that interaction class.}

\defentry{base object model (BOM)}{A piece-part of a conceptual model, simulation object model (SOM), or federation object model (FOM) that can be used as a building block in the development and/or extension of a simulation or federation.}

\defentry{best effort}{A transportation type that does not provide any guarantee regarding delivery (messages may be out of order, duplicate, or dropped).}

\defentry{candidate discovery class}{The registered class of an object instance, if subscribed. If the registered class of an object instance is not subscribed, the closest superclass of the registered class of the object instance to which the joined federate is subscribed. Candidate discovery class pertains to object instances only.}

\defentry{candidate received class}{The sent class of an interaction, if subscribed. If the sent class of an interaction is not subscribed, the closest superclass of the sent class of the interaction to which the joined federate is subscribed. Candidate received class pertains to interactions only.}

\defentry{class}{A description of a group of items with similar properties, common behavior, common relationships, and common semantics.}

\defentry{class attribute}{A named characteristic of an object class denoted by a pair composed of an object class designator and an attribute designator.}

\defentry{class hierarchy}{A specification of a class-subclass or ``is-a'' relationship between classes in a given\\ domain.}

\defentry{coadunate}{To attach an object instance handle and an object instance name to an object instance. The\\ object instance name can be provided specifically by the federate or implicitly by the runtime infrastructure (RTI).}

\defentry{collection}{A group of elements in which an element may occur multiple times (this definition corresponds to the mathematical notion of a bag).}

\defentry{collection of pairs}{A group of pairs in which multiple pairs may have the same first component and a given pair may occur multiple times.}

\defentry{compliant object model}{A High Level Architecture (HLA) federation object model (FOM) or simulation object model (SOM) that fully conforms with all of the rules and constraints specified in object model template (OMT).}

\newpage

\defentry{constrained set of pairs}{A group of pairs in which no two pairs have the same first component (this definition corresponds to the mathematical notion of a function). An example would be a group of instance attribute and value pairs; each instance attribute may have at most one value associated with it.}

\defentry{corresponding class attribute of an instance attribute}{The class attribute that, from the perspective of a given joined federate, is the class attribute of the joined federate's known class for the object instance containing the instance attribute that has the same attribute designator as the instance attribute.}

\defentry{corresponding instance attributes of a class attribute}{The instance attributes that, from the perspective of a given joined federate, are}
\begin{hlalist}
    \item Unowned instance attributes of object instances that have a known class at the joined federate equal to the object class of the class attribute and that have the same attribute designator as the class attribute, or
    \item Instance attributes owned by the joined federate that belong to object instances that have a known class at the owning federate equal to the object class of the class attribute and that have the same attribute designator as the class attribute.
\end{hlalist}

\defentry{current federation object model (FOM)}{The union of the FOM modules and one Management Object Model (MOM) and Initialization Module (MIM) that have been specified in the creation of the federation execution or by any federate that has joined the federation execution. The sum operation is carried out according to the rules as prescribed in IEEE Std 1516.2-2010. When all FOM modules have been provided, the current FOM will be equal to the FOM; and before this step has happened, the current FOM will be a true subset of the FOM.}

\defentry{current federation object model (FOM) Document Data (FDD)}{The data and information used to\\ configure the federation execution that is the union of the FOM modules and one Management Object Model (MOM) and Initialization Module (MIM) that have been specified in the creation of the federation execution or by any federate that has joined the federation execution. The sum operation is carried out according to the rules as prescribed in IEEE Std 1516.2-2010. When all FOM modules have been provided, the current FDD will be equal to the FDD; and before this step has happened, the current FDD will be a true\\ subset of the FDD.}

\defentry{datatype}{A representation convention for a data element establishing its format, resolution, cardinality, and ordinality.}

\defentry{declared attributes}{The set of class attributes of a particular object class that are listed in the federation object model (FOM) as being associated with that object class in the object class hierarchy tree.}

\defentry{declared parameters}{The set of parameters of a particular interaction class that are listed in the federation object model (FOM) as being associated with that interaction class in the interaction class hierarchy tree.}

\defentry{default range}{A range lower bound and a range upper bound, defined in the federation object model (FOM) Document Data (FDD) and specified in terms of [0, the dimension's upper bound), for a dimension.}

\defentry{default region}{A multidimensional region provided by the runtime infrastructure (RTI) that is composed of one range for each dimension found in the federation object model (FOM) Document Data (FDD). The bounds of each of these ranges are [0, the range's dimension's upper bound). There is no way for a federate to refer to the default region.}
\hlaNote{See 9.1 in IEEE Std 1516.1-2010.}

\newpage

\defentry{delete}{To invoke the \hlaService{Delete Object Instance} service to eliminate a particular object instance. \hlaSeeAlso{\hlaTermRef{remove}.}}
\hlaNote{See 6.10 in IEEE Std 1516.1-2010.}

\defentry{designator}{A generic view of arguments referenced in service descriptions. This view improves clarity in situations when arguments (mostly identifiers) to services have different views (implementations) due to a particular programming language or application programmer's interface (API).}

\defentry{dimension}{A named interval. The interval is defined by an ordered pair of values, the first being the dimension lower bound and the second being the dimension upper bound. A runtime infrastructure (RTI) provides a predefined interval, whose lower and upper bounds are fixed as [0, the dimension's upper bound) as specified in the federation object model (FOM) Document Data (FDD). This interval provides a single basis for communication of all dimension-related data with an RTI. All normalized intervals communicated to the RTI will be subsets of this interval.}
\hlaNote{See 9.1 of IEEE Std 1516.1-2010.}

\defentry{discover}{To receive an invocation of the \hlaService{Discover Object Instance \textdagger}\footnote{All RTI-initiated services are denoted with a \textdagger\ (printer's dagger) after the service name.} service for a particular object\\ instance.}
\hlaNote{See 6.5 of IEEE Std 1516.1-2010.}

\defentry{discovered class}{The class that was an object instance's candidate discovery class at a joined federate when that object instance was discovered by that joined federate. \hlaSeeAlso{\hlaTermRef{candidate discovery class}.}}
\hlaNote{See 5.1.2 in IEEE Std 1516.1-2010.}

\defentry{exception}{Notification of any irregularity that may occur during a service invocation.}

\defentry{fault}{An event that prevents the entire federation from executing in a High Level Architecture (HLA) compliant manner.}

\defentry{federate}{An application that may be or is currently coupled with other software applications under a federation object model (FOM) Document Data (FDD) and a runtime infrastructure (RTI). \hlaSeeAlso{\hlaTermRef{federate application}; \hlaTermRef{joined federate}.}}

\defentry{federate application}{An application that supports the High Level Architecture (HLA) interface to a\\ runtime infrastructure (RTI) and that is capable of joining a federation execution. A federate application may join the same federation execution multiple times or may join multiple federation executions. However, each time a federate application joins a federation execution, it is creating a new joined federate. \hlaSeeAlso{\hlaTermRef{joined federate}.}}

\defentry{federation}{A named set of federate applications and a common federation object model (FOM) that are used as a whole to achieve some specific objective.}

\defentry{federation execution}{The actual operation, over time, of a set of joined federates that are interconnected by a runtime infrastructure (RTI).}

\defentry{federation objectives}{The statement of the problem that is to be addressed by the establishment and execution of a federation.}

\newpage

\defentry{federation object model (FOM)}{A specification defining the information exchanged at runtime to achieve a given set of federation objectives. This information includes object classes, object class attributes, interaction classes, interaction parameters, and other relevant information. The FOM is specified to the runtime infrastructure (RTI) using one or more FOM modules. The RTI assembles a FOM using these FOM modules and one Management Object Model (MOM) and Initialization Module (MIM), which is provided automatically by the RTI or, optionally, provided to the RTI when the federation execution is created.}

\defentry{federation object model (FOM) Document Data (FDD)}{The data and information in a FOM that are used by the \hlaService{Create Federation Execution} service and successive \hlaService{Join Federation Execution} service invocations to configure the federation execution.}

\defentry{federation object model (FOM) module}{A partial FOM (containing some or all object model template (OMT) tables) that specifies a modular component of a FOM. A FOM module may contain classes not inherent to it but upon which the FOM module depends, i.e., superclasses to the modular components. These superclasses will be included in the FOM module as either complete or scaffolding definitions.}

\defentry{federation requirements}{A statement that identifies a federation characteristic, constraint, process, or product that is unambiguous and testable and that is necessary for a federation to be acceptable for its intended use.}

\defentry{federation scenario}{A set of initial conditions and timeline of significant events used within a federation execution to achieve federation objectives.}

\defentry{handle}{An identifier originated/created by the runtime infrastructure (RTI) that is unique to a federation execution.}

\defentry{High Level Architecture (HLA) time axis}{A totally ordered sequence of values in which each value\\ typically represents an HLA instant of time in the physical system being modeled. For any two points, T1 and T2, on the time axis, if T1 $<$ T2, T1 represents an instant of time that occurs before the instant represented by T2.}

\defentry{inherited attribute}{A class attribute of an object class that was declared in a superclass of that object class in the object class hierarchy tree defined in the federation object model (FOM).}

\defentry{inherited parameter}{An interaction parameter that was declared in a superclass of that interaction class in the interaction class hierarchy tree defined in the federation object model (FOM).}

\defentry{in scope}{Of or pertaining to an instance attribute of an object instance for which all of the following apply:}
\begin{hlalist}
    \item The object instance is known to the joined federate.
    \item The instance attribute is owned by another joined federate.
    \item The instance attribute's corresponding class attribute is a one of the following:
    \begin{enumerate}[label=\arabic*),leftmargin=2.2em,labelsep=0.55em,itemsep=1pt,topsep=1pt,parsep=0pt,partopsep=0pt]
  \item A subscribed attribute of the known class of the object instance, or
  \item A subscribed attribute of the known class of the object instance with regions, and the update region set of the instance attribute at the owning joined federate overlaps the subscription region set of the instance attribute's corresponding class attribute at the known class of the instance attribute at the subscribing joined federate.
    \end{enumerate}
\end{hlalist}
\vspace{-7.5pt}
\hlaNote{See 6.1 of IEEE Std 1516.1-2010.}

\defentry{instance attribute}{A named characteristic of an object instance denoted by a pair composed of the object instance designator and the attribute designator.}

\newpage

\defentry{interaction}{An explicit action taken by a federate that may have some effect or impact on another federate within a federation execution.}

\defentry{interaction class}{A template for a set of characteristics that is common to a group of interactions. These characteristics correspond to the parameters that individual federates may associate with interactions.}

\defentry{interaction parameters}{The information associated with an interaction that a federate potentially affected by the interaction may receive to calculate the effects of that interaction on its current state.}

\defentry{joined federate}{A member of a federation execution, actualized by a federate application invoking the \hlaService{Join Federation Execution} service as prescribed in IEEE Std 1516.1-2010. \hlaSeeAlso{\hlaTermRef{federate application}.}}

\defentry{known class}{}
\begin{hlalist}
    \item An object instance's registered class if the joined federate knows about the object instance as a result of having registered it, or
    \item An object instance's discovered class if the joined federate knows about the object instance as a result of having discovered it.
\end{hlalist}

\defentry{known object instance}{An object instance that a given joined federate has either registered or discovered and that the joined federate has not subsequently deleted (globally or locally) or been notified to remove. \hlaSeeAlso{\hlaTermRef{register}; \hlaTermRef{discover}; \hlaTermRef{delete}; \hlaTermRef{local delete}; \hlaTermRef{remove}.}}

\defentry{last known good logical timestamp}{The last timestamp to which a lost joined federate was successfully granted, as seen from the remaining High Level Architecture (HLA) compliant federation.}

\defentry{local delete}{To invoke the \hlaService{Local Delete Object Instance} service to inform the runtime infrastructure (RTI) that it is to treat the specified object instance as if the RTI had never notified the joined federate to discover the object instance (however, the object instance is not to be eliminated and may be rediscovered). \hlaSeeAlso{\hlaTermRef{delete}.}}
\hlaNote{See 6.12 of IEEE Std 1516.1-2010.}

\defentry{logical time}{A federate's current point on the High Level Architecture (HLA) time axis. Federates making use of the time management services follow restrictions on what timestamps can be sent in timestamp order (TSO) messages (relative to their logical time) to ensure that federates receiving those messages receive them in TSO.}

\defentry{lookahead}{A nonnegative value that establishes a lower value on the timestamps that can be sent in timestamp order (TSO) messages by time-regulating joined federates. Once established, a joined federate's lookahead value may be changed only by using the \hlaService{Modify Lookahead} service. Each time-regulating joined federate must provide a lookahead value when becoming time-regulating.}

\defentry{lost joined federate}{A previously joined federate that has been disconnected from the runtime infrastructure (RTI) as a result of a fault so that the joined federate can no longer continue in the federation execution in a High Level Architecture (HLA) compliant manner.}

\defentry{Management Object Model (MOM)}{A group of predefined High Level Architecture (HLA) constructs (object and interaction classes) that provide the following:}
\begin{hlalist}
    \item Access to federation execution operating information
    \item Insight into the operations of joined federates and the runtime infrastructure (RTI)
    \item Control of the functioning of the RTI, the federation execution, and the individual joined federates
\end{hlalist}
